% !TEX program = lualatex
% 自動生成: mar.report.latex — 手で編集した場合は再生成で上書きされる
\documentclass[paper=a4,fontsize=10.5pt,jafontsize=10.5pt]{jlreq}

\usepackage{luatexja}
\usepackage[haranoaji]{luatexja-preset}
\usepackage{amsmath,amssymb,amsthm,mathtools}
\usepackage{bm}
\usepackage{graphicx}
\usepackage{booktabs}
\usepackage{listings}
\usepackage{xcolor}
\usepackage[margin=25mm]{geometry}
\usepackage[hidelinks]{hyperref}
\usepackage{url}

\lstset{
  basicstyle=\ttfamily\footnotesize,
  breaklines=true,
  frame=single,
  showstringspaces=false,
  numbers=left,
  numberstyle=\tiny\color{gray},
  captionpos=b,
}

\theoremstyle{plain}
\newtheorem{theorem}{定理}[section]
\newtheorem{proposition}[theorem]{命題}
\newtheorem{lemma}[theorem]{補題}
\newtheorem{corollary}[theorem]{系}
\newtheorem{conjecture}[theorem]{予想}
\theoremstyle{definition}
\newtheorem{definition}[theorem]{定義}
\newtheorem{example}[theorem]{例}
\newtheorem{remark}[theorem]{注意}

\renewcommand{\proofname}{証明}

\title{葉数と局所独立数の平均: Written on the Wall II 予想 2 に対する二重星定理と局所版予想}
\author{自動研究パイプライン \texttt{mar}（Claude Opus 5 + 検証器）}
\date{2026年07月27日}

\begin{document}
\maketitle

\begin{abstract}
$G$ を位数 $n$ の連結グラフ、$L_s(G)$ を\textbf{葉数} (全域木がもつ葉の最大個数)、$\ell(v) = \alpha(G[N(v)])$ を\textbf{局所独立数}、$\overline{\ell}(G) = \frac{1}{n}\sum_v \ell(v)$ をその平均とする。Graffiti.pc の未解決予想 \emph{Written on the Wall II} 予想 2 (以下、予想 \ref{conj:wowii2}) は $2(\overline{\ell}(G) - 1) \le L_s(G)$ を主張する。本稿は 3 つの寄与を与える。第 1 に\textbf{二重星定理}: 連結グラフの\textbf{任意の辺} $uv$ に対し $L_s(G) \ge |N(u) \cup N(v)| - 2$ (定理 \ref{thm:doublestar})。証明は「部分木を全域木へ延長しても葉は減らない」という一行の補題だけを使う。第 2 に、この定理が予想 \ref{conj:wowii2} を\textbf{局所的で多項式時間に確かめられる}主張へ帰着させることを指摘し、その主張を予想 B' $\max_{uv \in E} |N(u) \cup N(v)| \ge 2\overline{\ell}(G)$ として提出する (予想 \ref{conj:bprime})。B' は予想 \ref{conj:wowii2} を含意し (系 \ref{cor:bprime})、三角形をもたないグラフでは Cauchy--Schwarz の 3 行で証明できる (定理 \ref{thm:trianglefree})。第 3 に、自明帯 $\overline{\ell} \le 2$ と $\Delta$ 帯 $2\overline{\ell} \le \Delta + 2$ で予想 \ref{conj:wowii2} を証明し (定理 \ref{thm:zones})、この 2 帯が実際に走査範囲の 99.656\% を覆うことを示す。連結グラフの完全リスト ($n \le 10$)、木 ($11 \le n \le 20$)、GENREG の連結正則グラフ、計 13,402,242 個について、各グラフに\textbf{葉集合 1 つ}を証人として付け、探索器とグラフ計算を共有しない検証器に (i) 補集合が連結支配集合であること、(ii) 予想 \ref{conj:wowii2}、(iii) 二重星定理、(iv) 予想 B' を再計算させた。走査した 13,402,242 個すべてで予想 \ref{conj:wowii2} は成立し、等号は 13 個で成立した。本稿が新たに提出する予想 \ref{conj:bprime} にも反例はなく、等号は 8,796 個で成立した。予想 \ref{conj:wowii2} 自体は\textbf{一般には未解決のままである}。
\end{abstract}

\noindent\textbf{キーワード:} graph theory、leaf number、connected domination、local independence、Graffiti、open problem、certificate


\section{はじめに}

$G$ を位数 $n = |V(G)|$、辺数 $m = |E(G)|$ の連結な単純グラフとする。
$N(v)$ を頂点 $v$ の\textbf{開}近傍、$\deg(v) = |N(v)|$、$\Delta(G)$ を
最大次数、$\alpha(\cdot)$ を独立数とする。
\[ \ell(v) \ = \ \alpha\bigl(G[N(v)]\bigr), \qquad
   \overline{\ell}(G) \ = \ \frac{1}{n} \sum_{v \in V(G)} \ell(v) \]
と置く。$\ell(v)$ は「$v$ の隣人のうち互いに隣接しないものを最大何個取れるか」
であり、$G$ が $v$ のまわりでどれだけ局所的に疎かを測る\textbf{局所独立数}で
ある。$L_s(G)$ は $G$ の全域木がもつ葉の最大個数 (\textbf{葉数}) を表す。

\begin{conjecture}[Graffiti.pc; WOWII Conjecture 2 \cite{wowii}]
\label{conj:wowii2}
連結グラフ $G$ に対し
\[ 2\bigl(\overline{\ell}(G) - 1\bigr) \ \le \ L_s(G) . \]
\end{conjecture}

$n$ 倍して分母を払うと、整数だけの同値な形
\[ n \, L_s(G) \ \ge \ 2 S(G) - 2n, \qquad
   S(G) = \sum_{v \in V(G)} \ell(v) \tag{$\star$} \]
になる。本稿の機械照合はこの $(\star)$ の形で行う (浮動小数点を使わない)。

この予想は DeLaViña の公開している未解決リスト \cite{wowii} に
2026 年 7 月 27 日に取得した時点でも状態 \texttt{O} (未解決) として載って
おり、同サイトの解決済みリストには現れない。左辺は多項式時間で計算できる
($\ell(v)$ は近傍という小さい部分グラフの独立数である) 一方、$L_s(G)$ の決定は
NP 困難である。予想 \ref{conj:wowii2} はしたがって「安価な量で計算困難な量を
下から押さえる」型の主張であり、Graffiti.pc が量産したこの型の予想のうち、
$L_s$ を扱うものの多くは DeLaViña--Waller \cite{dw2008} や
DeLaViña--Fajtlowicz--Waller \cite{dfw2005} で解決されたが、予想
\ref{conj:wowii2} は残っている。$L_s$ そのものの下界としては、次数条件に
基づく古典的なもの (立方体グラフに対する Griggs--Kleitman--Shastri
\cite{griggs} など) が知られているが、いずれも $\overline{\ell}$ という
局所量とは結びついていない。

\medskip
\noindent\textbf{本稿の寄与}。予想 \ref{conj:wowii2} は一般には証明できて
いない。得られたのは次の 3 つである。

\begin{enumerate}
\item \textbf{二重星定理} (定理 \ref{thm:doublestar}): 連結グラフの
  \textbf{すべての辺} $uv$ で $L_s(G) \ge |N(u) \cup N(v)| - 2$。
  三角形をもたないグラフで使われてきた二重星の議論を、
  $\deg u + \deg v$ ではなく\textbf{近傍の合併}で数えることによって、
  三角形の有無によらない形に持ち上げたものである。
\item \textbf{局所版予想 B'} (予想 \ref{conj:bprime}、本稿が提出):
  連結グラフは $\max_{uv \in E} |N(u) \cup N(v)| \ge 2\overline{\ell}(G)$
  を満たす。二重星定理と合わせると予想 \ref{conj:wowii2} が出る
  (系 \ref{cor:bprime})。B' には NP 困難な量が現れず、$O(nm)$ で判定できる。
  三角形をもたないグラフでは Cauchy--Schwarz で証明できる
  (定理 \ref{thm:trianglefree})。
\item \textbf{2 つの帯での証明と 13,402,242 個での機械照合}: 自明帯
  $\overline{\ell} \le 2$ と $\Delta$ 帯
  $2\overline{\ell} \le \Delta + 2$ で予想 \ref{conj:wowii2} を証明し
  (定理 \ref{thm:zones})、この 2 帯が走査範囲の 99.656\% を
  覆うことを示す。残る 46,096 個 (0.344\%) も、証人つきで
  機械照合した。
\end{enumerate}

\section{準備: 葉数と連結支配数}\label{sec:prelim}

$C \subseteq V(G)$ が\textbf{連結支配集合}であるとは、$C \ne \emptyset$、
$G[C]$ が連結、かつ $V \setminus C$ の各点が $C$ に隣接することをいう。

\begin{proposition}\label{prop:cds}
$G$ を連結グラフ、$L \subseteq V(G)$ とする。$V(G) \setminus L$ が連結支配
集合ならば $L_s(G) \ge |L|$ である。逆に $n \ge 3$ のとき、葉数を達成する
全域木の葉集合 $L$ に対し $V(G) \setminus L$ は連結支配集合であり、したがって
$L_s(G) = n - \gamma_c(G)$ ($\gamma_c$ は最小の連結支配集合の大きさ) が
成り立つ。
\end{proposition}

\begin{proof}
$C = V \setminus L$ が連結支配集合なら、$G[C]$ の全域木を 1 つ取り、$L$ の各点を
$C$ にある隣人 1 つへ結べば $G$ の全域木ができ、$L$ の点はすべて葉である。
逆に $T$ を全域木、$L$ をその葉集合とすると、木から葉を全部落としても連結性は
保たれるので $G[C] \supseteq T - L$ は連結であり、$L$ の各点は $T$ の中で
$C$ の点に隣接する。$n \ge 3$ では $C \ne \emptyset$ ($T$ の全点が葉なら
$T$ は $K_2$)。
\end{proof}

\begin{remark}
$n = 2$ は例外である。$K_2$ では両端が葉なので $L_s(K_2) = 2$ だが
$\gamma_c(K_2) = 1$ であり、$n - \gamma_c$ は 1 にしかならない。以下の議論と
機械照合はこの例外を明示的に扱う。なお $K_2$ では $\ell(v) = 1$ で
$S = 2$、$\overline{\ell} = 1$ だから $(\star)$ は $2 \cdot 2 \ge 0$ と
なって成立する (予想 \ref{conj:bprime} の方は $f(K_2) = 2 = 2\overline{\ell}$
で等号である)。
\end{remark}

命題 \ref{prop:cds} の前半だけを使う。すなわち\textbf{葉集合 1 つを見せれば
下界が立つ}。本稿の証明も機械照合も、この「証人」の形で進む。

\section{二重星定理}\label{sec:doublestar}

\begin{lemma}[部分木延長]\label{lem:extend}
$G$ を連結グラフ、$T_0$ を $G$ の部分グラフである木 ($|V(T_0)| \ge 2$) と
する。このとき $G$ は、葉を $T_0$ の葉と同じ個数以上もつ全域木を含む。
とくに $L_s(G) \ge (\text{$T_0$ の葉の個数})$。
\end{lemma}

\begin{proof}
$G$ は連結なので、$V(T_0) \ne V(G)$ である限り $V(T_0)$ に隣接する頂点 $w$ が
あり、その隣人 $p \in V(T_0)$ を 1 つ選んで辺 $pw$ を足せばまた木になる。この
操作で $w$ は新しい葉になり、葉でなくなり得るのは $p$ ただ 1 つである。よって
葉の個数は減らない。$V(T_0) = V(G)$ になるまで繰り返せばよい。
\end{proof}

\begin{theorem}[二重星定理]\label{thm:doublestar}
$G$ を連結グラフとする。$G$ の\textbf{任意の辺} $uv$ に対し
\[ L_s(G) \ \ge \ |N(u) \cup N(v)| - 2 . \]
\end{theorem}

\begin{proof}
$uv \in E$ より $u \in N(v)$、$v \in N(u)$ なので、
$A = (N(u) \cup N(v)) \setminus \{u, v\}$ と置くと
$|A| = |N(u) \cup N(v)| - 2$ である。$T_0$ を次の木とする: 頂点集合は
$\{u, v\} \cup A$、辺は $uv$、$u$ と $N(u) \setminus \{v\}$ の各点を
結ぶ辺、そして $v$ と $A \setminus N(u)$ の各点を結ぶ辺。$A$ の各点はこの
中で次数 1 なので、$T_0$ は葉を少なくとも $|A|$ 個もつ二重星である
($A = \emptyset$ なら $T_0 = K_2$ で主張は自明)。補題 \ref{lem:extend} を
$T_0$ に適用すればよい。
\end{proof}

証明は初等的だが、右辺が\textbf{辺ごとに}効くこと、そして
$|N(u) \cup N(v)|$ が $\deg u + \deg v$ ではなく\textbf{合併}であることが
要点である。三角形があると $\deg u + \deg v$ は共通近傍を二重に数えてしまい
偽の下界になるが、合併ならその心配がない。同じ頂点に星を立てるだけの議論と
比べると、$u$ と $v$ の近傍が食い違っている分だけ得をしている。

\begin{corollary}\label{cor:delta}
連結グラフ $G$ に対し $L_s(G) \ge \Delta(G)$。
\end{corollary}

\begin{proof}
最大次数の頂点 $v^*$ に星 $T_0$ ($v^*$ とその $\Delta$ 個の隣人) を立てると葉が
$\Delta$ 個ある ($\Delta = 1$ なら $G = K_2$ で $L_s = 2 > 1$)。補題
\ref{lem:extend} による。
\end{proof}

\section{局所版予想 B'}\label{sec:bprime}

定理 \ref{thm:doublestar} により、予想 \ref{conj:wowii2} は NP 困難な
$L_s$ を含まない主張へ帰着する。

\begin{conjecture}[本稿]\label{conj:bprime}
連結グラフ $G$ に対し
\[ f(G) \ := \ \max_{uv \in E(G)} |N(u) \cup N(v)|
   \ \ge \ 2\,\overline{\ell}(G) . \]
\end{conjecture}

\begin{corollary}\label{cor:bprime}
予想 \ref{conj:bprime} は予想 \ref{conj:wowii2} を含意する。
\end{corollary}

\begin{proof}
$f(G)$ を達成する辺に定理 \ref{thm:doublestar} を使うと
$L_s(G) \ge f(G) - 2 \ge 2\overline{\ell}(G) - 2$。
\end{proof}

B' の利点は 2 つある。第 1 に、$f(G)$ も $\overline{\ell}(G)$ も
$O(nm)$ 程度で計算できるので、反例探索が全域木の最適化から解放される
(実際、本稿の走査でも $L_s$ を厳密に解いたのは 13 回、
全体の 0.000097\% だけである)。第 2 に、主張が\textbf{辺 1 本の
近傍という局所的な対象}だけを見ているので、次数列や局所構造からの議論が
そのまま使える。実際、三角形がなければ 3 行で片づく。

\begin{theorem}\label{thm:trianglefree}
$G$ が三角形をもたない連結グラフならば予想 \ref{conj:bprime} が成り立つ。
したがって予想 \ref{conj:wowii2} も成り立つ。
\end{theorem}

\begin{proof}
三角形がないので各 $N(v)$ は独立集合であり $\ell(v) = \deg(v)$、すなわち
$2\overline{\ell}(G) = \frac{2}{n}\sum_v \deg(v) = \frac{4m}{n}$。
また隣接する $u, v$ が共通近傍をもてば三角形ができるので
$N(u) \cap N(v) = \emptyset$、ゆえに
$|N(u) \cup N(v)| = \deg u + \deg v$。Cauchy--Schwarz より
\[ \sum_{uv \in E} \bigl(\deg u + \deg v\bigr)
   \ = \ \sum_{v \in V} \deg(v)^2
   \ \ge \ \frac{1}{n}\Bigl(\sum_{v} \deg v\Bigr)^2
   \ = \ \frac{4m^2}{n} \]
なので、平均を上回る辺が存在する:
$f(G) \ge \frac{1}{m} \cdot \frac{4m^2}{n} = \frac{4m}{n}
 = 2\overline{\ell}(G)$。
\end{proof}

三角形をもたないグラフでの予想 \ref{conj:wowii2} は、連結支配数の上界を扱う
Mukwembi \cite{mukwembi} の議論と重なる可能性がある (本稿では原論文の
主張を突き合わせていない。第 \ref{sec:disc} 節を参照)。定理
\ref{thm:trianglefree} の意義は、証明が 3 行で済むことと、そこで実際に
示されているのが予想 \ref{conj:wowii2} より強い B' であることにある。

\section{2 つの帯と、そこで閉じない部分}\label{sec:zones}

三角形があるときは $\ell(v) < \deg(v)$ になり得て上の計算が崩れる。しかし
$\overline{\ell}$ が小さければ主張自体が弱くなる。

\begin{theorem}\label{thm:zones}
連結グラフ $G$ が次のいずれかを満たせば予想 \ref{conj:wowii2} が成り立つ。
\begin{enumerate}
\item (自明帯) $\overline{\ell}(G) \le 2$、すなわち $2S \le 4n$。
\item ($\Delta$ 帯) $2\overline{\ell}(G) \le \Delta(G) + 2$、すなわち
  $2S \le n(\Delta + 2)$。
\end{enumerate}
\end{theorem}

\begin{proof}
(1) 連結グラフの全域木は必ず 2 枚以上の葉をもつので $L_s(G) \ge 2 \ge
2(\overline{\ell} - 1)$。(2) 系 \ref{cor:delta} より
$L_s(G) \ge \Delta \ge 2\overline{\ell} - 2$。
\end{proof}

木はつねに自明帯にある: 木では $\ell(v) = \deg(v)$ なので
$S = 2(n-1) < 2n$。残るのは
\[ 2S \ > \ n\bigl(\Delta(G) + 2\bigr) \tag{H} \]
を満たすグラフ、すなわち「平均の局所独立数が最大次数に比べて大きい」層である。
$\ell(v) \le \deg(v) \le \Delta$ なので (H) は $\overline{\ell}$ が
$\Delta$ に近いことを要求し、そのようなグラフは実際には少ない
(表 \ref{tab:zones})。

\begin{table}[t]
\centering
\caption{連結グラフの完全リストにおける帯の分布。「残り」が (H) を満たす層で、
本稿が定理として閉じられていない部分である。}\label{tab:zones}
\begin{tabular}{rrrrrr}
\toprule
$n$ & 個数 & 自明帯 & $\Delta$ 帯 & 残り & 残りの割合 (\%) \\
\midrule
2 & 1 & 1 & 0 & 0 & 0.000 \\
3 & 2 & 2 & 0 & 0 & 0.000 \\
4 & 6 & 6 & 0 & 0 & 0.000 \\
5 & 21 & 20 & 1 & 0 & 0.000 \\
6 & 112 & 96 & 14 & 2 & 1.786 \\
7 & 853 & 570 & 276 & 7 & 0.821 \\
8 & 11,117 & 4,587 & 6,507 & 23 & 0.207 \\
9 & 261,080 & 45,253 & 215,593 & 234 & 0.090 \\
10 & 11,716,571 & 523,962 & 11,190,131 & 2,478 & 0.021 \\
\bottomrule
\end{tabular}
\end{table}

走査した 13,402,242 個のうち、自明帯が 1,920,364 個、$\Delta$ 帯が
11,435,782 個で、合わせて 13,356,146 個 (99.656\%) が定理
\ref{thm:zones} で閉じる。残りは 46,096 個 (0.344\%) で、
これらについては証人つきの機械照合しか持っていない。

\section{機械照合}\label{sec:check}

\subsection*{走査した族}

連結グラフの完全リスト ($2 \le n \le 10$、McKay \cite{mckay})、木
($ 11 \le n \le 20$、同)、GENREG \cite{genreg} の連結正則グラフ
(12, 3), (14, 3), (16, 3), (18, 3), (11, 4), (12, 4), (13, 4), (12, 5), (11, 6) を走査した。合計 13,402,242 個 (うち木 1,345,823 個、
正則グラフ 66,656 個)。族の個数は検証器が自分の持つ公表値と突き合わせて
おり、族を落とした証明書は検証に通らない。

\subsection*{証人}

各グラフに\textbf{葉集合 1 つ}を証人として付ける (32 bit のビットマスク 1 個、
族ごとに gzip でまとめ、SHA-256 を証明書に記録)。証人は 3 段で作る。

\begin{enumerate}
\item まず定理 \ref{thm:doublestar} の証明そのもの、すなわち $f(G)$ を
  達成する辺の二重星を延長して得られる葉集合 (13,402,070 個、
  99.999\% はこれで $(\star)$ が閉じた)。
\item 足りなければ貪欲な連結支配集合の補集合 (159 個)。
\item それでも足りなければ葉数を厳密に最大化する (13 個、
  厳密計算の呼び出しは 13 回)。
\end{enumerate}

\subsection*{検証器が再計算するもの}

検証器 (\texttt{mar.checkgraph}) は探索器とグラフ計算のコードを共有せず、
グラフを集合表現で持ち直したうえで、各グラフについて次を独立に計算する
(ただし下界の閾値 $2S - 2n$ と帯の判定式そのものは主張の\emph{定義}で
あって計算ではないので、この 2 つだけは探索器と共通の実装を使う)。

\begin{itemize}
\item $\ell(v) = \alpha(G[N(v)])$ と $S(G) = \sum_v \ell(v)$、および帯の判定。
\item 証人 $L$ について、$V \setminus L$ が空でなく、連結で、支配的であること
  (命題 \ref{prop:cds}) --- これが破れていれば $L$ は葉集合ではない。
\item $(\star)$: $n|L| \ge 2S - 2n$。
\item 定理 \ref{thm:doublestar}: $|L| \ge f(G) - 2$ ($f$ も検証器が
  別実装で計算する)。
\item 予想 \ref{conj:bprime}: $n f(G) - 2S \ge 0$ と、その最小値・等号個数。
\item 等号が記録されたグラフについては、葉数が $|L| + 1$ 以上にならないことを
  厳密に解き直す。逆に等号リストに載っていないグラフは $(\star)$ が
  \textbf{狭義に}閉じることを要求する (等号を隠せない)。
\end{itemize}

\subsection*{結果}

走査した 13,402,242 個すべてで予想 \ref{conj:wowii2} は成立し、等号は 13 個で成立した。本稿が新たに提出する予想 \ref{conj:bprime} にも反例はなく、等号は 8,796 個で成立した。

\begin{table}[t]
\centering
\caption{族ごとの内訳。「二重星」「貪欲」「厳密」は証人がどの段で得られたか
(厳密は探索器の自己申告で、検証器は再計算しない)。「等号」は
予想 \ref{conj:wowii2} の等号、「B' 等号」は予想 \ref{conj:bprime} の
等号の個数。}\label{tab:fams}
\begin{tabular}{lrrrrrrr}
\toprule
種別 & $n$ & 個数 & 二重星 & 貪欲 & 厳密 & 等号 & B' 等号 \\
\midrule
連結グラフ & 2 & 1 & 1 & 0 & 0 & 0 & 1 \\
連結グラフ & 3 & 2 & 2 & 0 & 0 & 0 & 0 \\
連結グラフ & 4 & 6 & 5 & 0 & 1 & 1 & 1 \\
連結グラフ & 5 & 21 & 20 & 0 & 1 & 1 & 1 \\
連結グラフ & 6 & 112 & 110 & 0 & 2 & 2 & 2 \\
連結グラフ & 7 & 853 & 852 & 0 & 1 & 1 & 1 \\
連結グラフ & 8 & 11,117 & 11,113 & 1 & 3 & 3 & 4 \\
連結グラフ & 9 & 261,080 & 261,079 & 0 & 1 & 1 & 1 \\
連結グラフ & 10 & 11,716,571 & 11,716,567 & 1 & 3 & 3 & 10 \\
木 & 11 & 235 & 235 & 0 & 0 & 0 & 0 \\
木 & 12 & 551 & 551 & 0 & 0 & 0 & 0 \\
木 & 13 & 1,301 & 1,301 & 0 & 0 & 0 & 0 \\
木 & 14 & 3,159 & 3,159 & 0 & 0 & 0 & 0 \\
木 & 15 & 7,741 & 7,741 & 0 & 0 & 0 & 0 \\
木 & 16 & 19,320 & 19,320 & 0 & 0 & 0 & 0 \\
木 & 17 & 48,629 & 48,629 & 0 & 0 & 0 & 0 \\
木 & 18 & 123,867 & 123,867 & 0 & 0 & 0 & 0 \\
木 & 19 & 317,955 & 317,955 & 0 & 0 & 0 & 0 \\
木 & 20 & 823,065 & 823,065 & 0 & 0 & 0 & 0 \\
3-正則 & 12 & 85 & 77 & 8 & 0 & 0 & 22 \\
3-正則 & 14 & 509 & 493 & 16 & 0 & 0 & 110 \\
3-正則 & 16 & 4,060 & 4,022 & 38 & 0 & 0 & 792 \\
3-正則 & 18 & 41,301 & 41,218 & 83 & 0 & 0 & 7,805 \\
4-正則 & 11 & 265 & 264 & 1 & 0 & 0 & 2 \\
4-正則 & 12 & 1,544 & 1,539 & 5 & 0 & 0 & 12 \\
4-正則 & 13 & 10,778 & 10,772 & 6 & 0 & 0 & 31 \\
5-正則 & 12 & 7,848 & 7,847 & 0 & 1 & 1 & 1 \\
6-正則 & 11 & 266 & 266 & 0 & 0 & 0 & 0 \\
\midrule
合計 & & 13,402,242 & 13,402,070 & 159 &
13 & 13 & 8,796 \\
\bottomrule
\end{tabular}
\end{table}

\subsection*{等号の分類}

等号成立グラフは 13 個すべてを証明書に列挙してあり、検証器はその各々について葉数を厳密に解き直して「これ以上葉を増やせない」ことを確かめている。表 \ref{tab:eq} がその全体である。走査範囲で等号が成立するのは
13 個で、そのうち 7 個は閉路 $C_n$ である
($\Delta = 2$ かつ $S = 2n$ の連結グラフは閉路に限る)。閉路では
$\ell(v) = 2$ が全頂点で成り立って $\overline{\ell} = 2$ となり、
全域木は道なので $L_s = 2 = 2(\overline{\ell} - 1)$ となる。閉路以外で同定できたのは $K_{3,3}$ ($n = 6$, $r = 3$, $L_s = 4$)、$Q_3$ ($n = 8$, $r = 3$, $L_s = 4$)、$K_{4,4}$ ($n = 8$, $r = 4$, $L_s = 6$)、$K_{5,5} - M$ ($n = 10$, $r = 4$, $L_s = 6$)、$K_{5,5}$ ($n = 10$, $r = 5$, $L_s = 8$)、$K_{6,6} - M$ ($n = 12$, $r = 5$, $L_s = 8$) である。ここで $M$ は完全マッチングであり、$K_{k,k} - M$ は冠グラフ (crown graph) である。いずれも $S = n\Delta$、すなわち $\ell \equiv \deg \equiv \Delta$ なので、三角形をもたない $\Delta$-正則グラフであり、等号は $L_s = 2\Delta - 2$ ちょうどであることを意味する。

\begin{table}[t]
\centering
\caption{予想 \ref{conj:wowii2} の等号が成立するグラフ。$f$ は
$\max_{uv \in E}|N(u) \cup N(v)|$。名前の列の「---」は本稿で同定して
いないことを表す。}\label{tab:eq}
\begin{tabular}{llrrrrrr}
\toprule
名前 & graph6 & $n$ & $L_s$ & $S$ & $\overline{\ell}$ & $\Delta$ & $f$ \\
\midrule
$C_{4}$ & \texttt{C]} & 4 & 2 & 8 & 2.000 & 2 & 4 \\
$C_{5}$ & \texttt{DUW} & 5 & 2 & 10 & 2.000 & 2 & 4 \\
$C_{6}$ & \texttt{EEh\_} & 6 & 2 & 12 & 2.000 & 2 & 4 \\
$K_{3,3}$ & \texttt{EFz\_} & 6 & 4 & 18 & 3.000 & 3 & 6 \\
$C_{7}$ & \texttt{FCp`\_} & 7 & 2 & 14 & 2.000 & 2 & 4 \\
$C_{8}$ & \texttt{G?qa`\_} & 8 & 2 & 16 & 2.000 & 2 & 4 \\
$Q_3$ & \texttt{G?zTb\_} & 8 & 4 & 24 & 3.000 & 3 & 6 \\
$K_{4,4}$ & \texttt{G?\textasciitilde{}vf\_} & 8 & 6 & 32 & 4.000 & 4 & 8 \\
$C_{9}$ & \texttt{H?bB@\_W} & 9 & 2 & 18 & 2.000 & 2 & 4 \\
$C_{10}$ & \texttt{I?BDA\_gE?} & 10 & 2 & 20 & 2.000 & 2 & 4 \\
$K_{5,5} - M$ & \texttt{I?BvUqw]?} & 10 & 6 & 40 & 4.000 & 4 & 8 \\
$K_{5,5}$ & \texttt{I?B\textasciitilde{}vrw\}?} & 10 & 8 & 50 & 5.000 & 5 & 10 \\
$K_{6,6} - M$ & \texttt{KsaBrhkV@w?\textasciicircum{}} & 12 & 8 & 60 & 5.000 & 5 & 10 \\
\bottomrule
\end{tabular}
\end{table}

等号が各位数で数個しか起きず、しかもそのすべてが名前のついた高度に対称なグラフであることは、予想 \ref{conj:wowii2} の下界が一般には
\textbf{大きく緩い}ことを示している。実際、表 \ref{tab:fams} のとおり
大多数は二重星だけで閉じる。

\section{議論}\label{sec:disc}

\subsection*{何が新しいか}

\begin{enumerate}
\item 定理 \ref{thm:doublestar} は、二重星の議論を三角形の有無によらない形
  ($\deg u + \deg v$ ではなく $|N(u) \cup N(v)|$) に整えたものである。
  証明は補題 \ref{lem:extend} 1 つで済み、$L_s$ の下界としてすぐ使える。
\item 予想 \ref{conj:bprime} は本稿が提出する新しい主張で、予想
  \ref{conj:wowii2} を含意しながら NP 困難な量を含まない。予想
  \ref{conj:wowii2} への攻撃をこの局所的な不等式へ移せることが、本稿の
  主要な観察である。
\item 定理 \ref{thm:trianglefree} は三角形をもたない場合の 3 行の証明を
  与える。しかも証明されるのは B' であって、予想 \ref{conj:wowii2} より
  強い。
\end{enumerate}

\subsection*{限界}

\textbf{予想 \ref{conj:wowii2} は一般には未解決のままである}。閉じて
いないのは (H) を満たす層、すなわち三角形をもち、かつ平均の局所独立数が最大
次数に近いグラフである。走査範囲ではこの層は 46,096 個
(0.344\%) しかないが、$n$ を大きくしたときにこの割合がどう振る舞う
かは本稿では調べていない。有限範囲の照合はいくら広くても証明ではない。

先行研究の網羅的な確認も行っていない。とくに Mukwembi \cite{mukwembi} の
定理の正確な主張は原論文で突き合わせておらず、定理
\ref{thm:trianglefree} が扱う三角形なしの場合が既知である可能性がある。
本文でその可能性を明示したのはこのためである。予想
\ref{conj:bprime} についても、同値な主張が別の言葉で既知でないことを
確かめてはいない。

機械照合の側にも境界がある。検証器が再計算するのは本節に挙げた量だけで、
探索器が自己申告する費用 (厳密計算の回数、証人がどの段で得られたか) は検証の
対象外である。また証人は「葉数の下界」しか与えないので、等号として記録された
グラフ以外について $L_s$ の真の値は確かめていない (必要がない)。

\subsection*{次に何をすべきか}

予想 \ref{conj:bprime} を一般に証明することが、そのまま予想
\ref{conj:wowii2} の解決になる。三角形なしの証明が使った等式
$|N(u) \cup N(v)| = \deg u + \deg v$ は、三角形があると
$\deg u + \deg v - |N(u) \cap N(v)|$ に落ちる。一方で三角形は $\ell(v)$ を
$\deg(v)$ より小さくする方向に働くので、$|N(u) \cap N(v)|$ の損失と
$\deg(v) - \ell(v)$ の得を突き合わせる不等式が作れれば、Cauchy--Schwarz の
議論をそのまま延長できる可能性がある。表 \ref{tab:eq} の等号グラフが
すべて三角形をもたないことは、その方向を支持している。

\medskip
\noindent
機械照合は \texttt{python -m mar verify p0009\_wowii2\_leaf\_local\_indep} で再実行できる。


\section*{再現性}
本稿の主張はすべて、探索器とは独立に実装された検証器によって再検査されている。次のコマンドで再現できる。

\begin{center}\texttt{python -m mar verify p0009\_wowii2\_leaf\_local\_indep}\end{center}

\begin{center}\begin{tabular}{ll}\toprule
証明書ダイジェスト & \texttt{f7ad2d8e997f1dc4} \\
生成日時 (UTC) & 2026-07-27T03:14:36+00:00 \\
Python & 3.10.11 \\
実行環境 & Windows 11 (build 26200) \\
リポジトリ版 & \texttt{4239d05} \\
乱数種 & 0 \\
探索時間 & 1468.6 秒 \\
\bottomrule\end{tabular}\end{center}


\begin{thebibliography}{99}
\bibitem{wowii} E. DeLaViña, Written on the Wall II: Conjectures of Graffiti.pc, University of Houston--Downtown (Conjecture 2; 2026-07-27 取得). \url{http://cms.dt.uh.edu/faculty/delavinae/research/wowII/}
\bibitem{mukwembi} S. Mukwembi, Size, order, and connected domination, Canad. Math. Bull. 57 (2014) 141--144.
\bibitem{dw2008} E. DeLaViña, W. Waller, Spanning trees with many leaves and average distance, Electron. J. Combin. 15 (2008) \#R33.
\bibitem{dfw2005} E. DeLaViña, S. Fajtlowicz, W. Waller, Spanning trees with many leaves, DIMACS Ser. Discrete Math. Theoret. Comput. Sci. 69 (2005) 119--125.
\bibitem{griggs} J. R. Griggs, D. J. Kleitman, A. Shastri, Spanning trees with many leaves in cubic graphs, J. Graph Theory 13 (1989) 669--695.
\bibitem{mckay} B. D. McKay, A. Piperno, Practical graph isomorphism II, J. Symbolic Comput. 60 (2014) 94--112. データ: Combinatorial Data. \url{https://users.cecs.anu.edu.au/~bdm/data/graphs.html}
\bibitem{genreg} M. Meringer, Fast generation of regular graphs and construction of cages, J. Graph Theory 30 (1999) 137--146. \url{https://www.mathe2.uni-bayreuth.de/markus/reggraphs.html}
\end{thebibliography}

\end{document}
